% Part A: the whole project. Shared, identical in all three documents.
\section{The whole project}

\subsection{The problem}
A normal security scanner gives a fixed severity to each kind of finding. A
missing \texttt{X-Content-Type-Options} header is ``medium'' wherever it is
found: on a public brochure page, and on a login API that returns credentials.
The person reading the report then has to work out, by hand, which of a hundred
``medium'' findings actually matters. That triage is slow and easy to get
wrong. Tools such as OWASP ZAP and Nuclei do find real issues, but they rate
the issue, not the place it sits in.\src{docs/supervisor/scanners.md}

\subsection{The idea}
CAWASMA scores the \emph{same} finding differently according to where it occurs
and what the endpoint returned. It does this with the standard CVSS v3.1
\emph{environmental} metrics, so the result is not an invented number: it is a
recomputation of the official score with a vector anyone can check on the FIRST
calculator.

\keybox{\textbf{One-paragraph pitch.} CAWASMA is a defensive, black-box,
HTTPS-only scanner for web apps and REST APIs. Each of its checks carries a
fixed CVSS v3.1 base vector. For every finding it works out how sensitive the
endpoint is (from the URL path) and what sensitive data the response holds
(credentials, personal data, financial data, internal IPs). It turns those into
CVSS environmental security requirements (CR/IR/AR) and recomputes the score, so
a login API and a public page that share a weakness no longer share a priority.
The point of the project is \emph{prioritisation, not detection}.}

\subsection{Architecture}
A scan is a pipeline. A URL goes in; a ranked, explained report comes out.

\begin{center}
\resizebox{0.98\linewidth}{!}{%
\begin{tikzpicture}[
  font=\footnotesize,
  >={Stealth[length=2mm]},
  box/.style={draw=rule,fill=soft,rounded corners=2pt,minimum height=8mm,inner sep=3pt,text width=22mm,align=center},
  core/.style={box,fill=brand!12,draw=brand},
  node distance=4mm and 6mm,
]
\node[box] (url) {Target URL};
\node[core,right=of url] (eng) {HTTP engine\\\textit{SSRF guard}};
\node[core,right=of eng] (crawl) {Crawler\\\textit{find endpoints}};
\node[core,right=of crawl] (checks) {Checks\\\textit{58 detectors}};
\node[core,below=8mm of checks] (sens) {Sensitivity\\\textit{class of path}};
\node[core,left=of sens] (body) {Body signals\\\textit{data in reply}};
\node[core,left=of body] (cvss) {\texttt{cvss.py}\\\textit{environmental}};
\node[box,left=of cvss] (db) {Database};
\node[box,below=8mm of db] (api) {REST API\\+ SSE};
\node[box,right=of api] (ui) {Dashboard};
\node[box,right=of ui] (rep) {Reports\\\textit{JSON/PDF}};
\draw[->] (url)--(eng); \draw[->] (eng)--(crawl); \draw[->] (crawl)--(checks);
\draw[->] (checks)--(sens); \draw[->] (sens)--(body); \draw[->] (body)--(cvss);
\draw[->] (cvss)--(db); \draw[->] (db)--(api); \draw[->] (api)--(ui); \draw[->] (ui)--(rep);
\draw[->,dashed] (checks.north) to[out=90,in=90] node[above,font=\scriptsize]{findings} (cvss.north);
\end{tikzpicture}}
\end{center}

\noindent The stages are: the \textbf{engine} fetches safely (one SSRF check per
redirect hop); the \textbf{crawler} discovers endpoints (seed, HTML, JavaScript,
a wordlist, and OpenAPI paths); the \textbf{checks} look for misconfigurations;
the \textbf{sensitivity} classifier labels each endpoint CRITICAL/HIGH/MEDIUM/LOW
from its path; the \textbf{body classifier} looks for sensitive data in the
reply; \texttt{cvss.py} turns those two into CR/IR/AR and recomputes the score;
the result is stored, served by the API over REST and Server-Sent Events, shown
on the dashboard, and exported as reports.

\subsection{Technology, and why}
\begin{center}
\small
\begin{tabularx}{\linewidth}{@{}p{34mm}X@{}}
\toprule
\textbf{Part} & \textbf{Choice and reason}\\
\midrule
Scanner core & Python 3.12 with \texttt{httpx} + \texttt{asyncio}: endpoints are analysed concurrently, and the core stays database-free so it is easy to test.\\
Scoring & The \texttt{cvss} library, so the environmental maths is the published CVSS v3.1 formula, not ours.\\
NLP & Sentence-Transformers (\texttt{all-MiniLM-L6-v2}) and spaCy for the sensitivity class; the engine still works with the built-in classifier when these are absent.\\
API & FastAPI, with Pydantic schemas and a frozen contract; Celery + Redis run scans in the background with live status.\\
Database & PostgreSQL in production, SQLite in tests; the models stay portable with a \texttt{JSONType} column and database-level cascades.\\
TLS & SSLyze (with a plain-\texttt{ssl} fallback) for the transport checks.\\
Reports & One HTML template through WeasyPrint for the PDF, so HTML and PDF never drift.\\
Frontend & Next.js 16 (App Router), React 19, TypeScript, Tailwind CSS v4, shadcn/ui, TanStack Query.\\
Tests & \texttt{pytest} (backend), Vitest + Testing Library and Playwright (frontend).\\
\bottomrule
\end{tabularx}
\end{center}

\subsection{Team roles and the shared contracts}
Three people own three areas, and they meet at three written contracts, frozen
so one person's change cannot silently break another's code.\src{CLAUDE.md; .github/CODEOWNERS}

\begin{center}
\small
\begin{tabularx}{\linewidth}{@{}p{24mm}X@{}}
\toprule
\textbf{Member} & \textbf{Owns}\\
\midrule
Asim (22P-9229) & Check specifications and their CVSS v3.1 vectors, the environmental mapping, chain templates, the sensitivity and response (ML/NLP) classifiers, and the whole evaluation lab.\\
Neelam (22P-9353) & The Next.js web app, the \texttt{frontend} service and its CI job.\\
Atta (22P-9236) & The backend: HTTP engine, crawler, check functions, CVSS scoring code, attack-chain engine, reports, FastAPI app, Celery workers, database and migrations, plus the shared infrastructure and reviews.\\
\bottomrule
\end{tabularx}
\end{center}

\begin{description}[leftmargin=0pt,itemsep=3pt]
\item[API contract (v1).] \texttt{app/api/schemas/contract.py} and
\texttt{docs/api-contract.md}: the field names the frontend reads
(\texttt{check\_id}, \texttt{base\_score}, \texttt{final\_score}, status
\texttt{DONE}).
\item[ML / CVSS contract.] \texttt{app/core/contracts.py} and
\texttt{docs/cvss-mapping-contract.md}: the types the engine and Asim's
classifiers exchange, so the engine keeps working with the built-in
classifiers when the ML packages are missing.
\item[Check-spec format.] \texttt{app/core/specs.py} and
\texttt{docs/check-spec-format.md}: how a check is declared so the registry can
run it.
\end{description}

\paragraph{Two scoring states.} Not every check is re-scored the new way yet.
Of 62 check ids, 25 have an approved specification and are scored through the
CVSS v3.1 environmental metrics; the other 37 still use an interim formula,
$\min(10,\ \text{base}\times\text{context multiplier}+\text{response bonus})$,
and each finding records which method scored it
(\texttt{scoring\_method}).\src{make coverage-stats; docs/06-scoring-model.md}
